% 397_Bell_Synchronisation_En_ch.tex
% Johann Pascher, 18 Nov. 2026
% Bell violation through synchronization: pilot-wave droplets and the deterministic reading of the FFGFT

\providecommand{\BM}{\textbf{[B]}}
\providecommand{\KM}{\textbf{[K]}}
\providecommand{\SM}{\textbf{[S]}}
\providecommand{\XM}{\textbf{[X]}}
\providecommand{\QM}{\textbf{[Q]}}
\providecommand{\EM}{\textbf{[E]}}

\chapter*{Bell Violation through Synchronization}
\addcontentsline{toc}{chapter}{Bell Violation through Synchronization}

\begin{abstract}
Two papers from pilot-wave hydrodynamics show that two classical droplets coupled through a common
wave field can violate the Bell inequality in its CHSH form: in a static test with $S=2.49$, and in a
reduced model with $S\approx3.15$ in the static case and up to $S\approx2.85$ in a test in which the
coupling is switched off before the measurement. The authors trace the violation to deterministic,
wave-mediated synchronization and to the dependence of the wave field on the measurement settings.
This document presents the papers with verbatim quotations, recomputes the bounds and places them in
the deterministic reading of the FFGFT (Docs.~230, 328, 334). The result has two sides. The droplets
show that a deterministic single run with contextual fixing of the outcomes, as Doc.~230 assumes,
can carry Bell violations. They also show that measurement dependence alone knows no Tsirelson
bound: both values of the reduced model lie above $2\sqrt2$, the static one by more than 200
standard errors. In the FFGFT the quantum bound follows from the geometric quotients of Doc.~230, not
from synchronization; dynamic synchronization is ruled out as the mechanism of the FFGFT correlation
because the coupling $K_{\mathrm{eff}}=2\pi\xi$ is about 6000 times weaker than in the droplet model.
\end{abstract}

% ============================================================
\section*{Status markers}
% ============================================================
\begin{center}
\begin{tabular}{@{}lp{11cm}@{}}
\textbf{[B]} & Algebraically proved. \\[2pt]
\textbf{[K]} & Numerically confirmed (check script, comparison with measured or literature values). \\[2pt]
\textbf{[S]} & Assumption, open, or classification without a calculation of its own. \\[2pt]
\textbf{[X]} & Excluded. \\[2pt]
\textbf{[Q]} & Taken from the literature, with source. \\[2pt]
\textbf{[E]} & Established physics, not derived in this document. \\
\end{tabular}
\end{center}
Check script: \texttt{2/python/Dok397\_Skripte/pruef\_397\_bell\_synchronisation.py}.
Verbatim quotations appear in the original language in quotation marks and are marked \QM.

% ============================================================
\section{Occasion and question}
% ============================================================

The FFGFT reads quantum mechanics deterministically. Every measurement is a single process whose
outcome is fixed by the position of the system and by a quantity $\lambda$ that is set in the
interaction with the measuring arrangement; the Born value $|\alpha|^2$ is a geometric quotient and
fundamentally not a probability (Doc.~230). Bell correlations are taken to be real, but as a
topological connection on $T^4$, not as action at a distance in open space (Docs.~230, 193). The
FFGFT thus belongs to the approaches that do not share one of the premises of Bell's theorem: fixing
the outcome in the interaction with the measuring arrangement makes the model contextual.

In August 2026 a paper appeared that works through exactly this route in a fully classical system
\QM~\cite{Lopez2026}. Two droplets walking on a vibrating oil bath and coupled through their common
wave field violate the CHSH inequality there, even after the coupling has been switched off before the
measurement. The question of this document is what that means for the deterministic reading of the
FFGFT: what it supports, what it does not, and what the FFGFT can learn from it.

% ============================================================
\section{The system: droplets riding their own wave}
% ============================================================

In 2005 Couder and Fort found that a millimetre-sized oil droplet on a vertically vibrating bath of
the same oil does not coalesce but bounces, exciting a standing Faraday wave at every impact
\QM~\cite{Couder2005}. Above a threshold it starts to move sideways, driven by the slope of its own
wave. The droplet is a particle, the wave a real field, and both are coupled through the sequence of
impacts. Bush's group has turned this into a field of its own, pilot-wave hydrodynamics, with a
growing list of classical counterparts to quantum effects: diffraction at single and double slits,
quantized orbits, tunnelling, statistics in confined corrals that follow the eigenmodes of the corral
\QM~\cite{BushOza2021}.

The decisive feature of the system is its memory. The wave decays slowly, and at every point it
carries the trace of all earlier impacts; the force on the droplet therefore depends on its entire
history. The quantum-like effects appear in the long-memory regime, when the waves are most
persistent \QM~\cite{Lopez2026}. The system is a macroscopic realization of de Broglie's double
solution: a real guiding wave and, alongside it, a statistical wave that emerges from many runs.

% ============================================================
\section{Two Bell tests with droplets}
% ============================================================

\subsection{The static test of Papatryfonos et al.}

Papatryfonos, Vervoort, Nachbin, Labousse and Bush simulated two droplets numerically, each in a pair
of wells between which it can tunnel; the two subsystems are joined by a coupling cavity that the
droplets do not enter but through which waves pass \QM~\cite{Papatryfonos2023}. The two-valued
outcome is the well in which the droplet sits at the time of measurement (outer $+1$, inner $-1$);
the measurement setting is the depth of the barrier between the wells. The authors themselves name
the key step: \enquote{the system geometry may serve as a proxy for analyzer settings}
\QM~\cite{Papatryfonos2023}. With symmetric settings they found a largest value of
$S_{\max}=2.49\pm0.04$, with a statistical confidence above 20 standard deviations, though only in a
narrow range of barrier depths \QM. They attribute the violation to the wave-mediated coupling: the
form of the pilot wave, and with it the position of the droplets, depends on the geometry of both
subsystems. A dynamic test, in which the coupling is switched off before the settings are changed,
was still open at the time.

\subsection{The reduced model of López, Valani, Li and Bush}

The August 2026 paper replaces the full wave equation with a sinusoidal wave kernel $W(x)=\cos x$,
which permits an exact reduction to ordinary differential equations \QM~\cite{Lopez2026}. Each droplet
sits in a double well; its motion follows from inertia, drag, the force of its own past wave, the
force of the partner's wave, and the external well. With memory variables $Y_{ij}$, $Z_{ij}$ that
carry the exponentially weighted history, a twelve-dimensional system of two coupled Lorenz-like
subsystems results. A switch $\zeta\in\{0,1\}$ turns the coupling between the droplets on or off.

The symmetry of the system provides two invariant synchronization manifolds, one for in-phase and one
for anti-phase motion. On them the droplets run periodically or chaotically in step; in the chaotic
regime the trajectory switches irregularly between the two manifolds. Depending on the well, that is,
depending on the setting, the trajectory stays longer near one or the other, and this produces the
correlations. In the static test the mean over 1000 random initial conditions converges to
$\langle S\rangle\approx3.15$ with a standard error of about $1.5\times10^{-3}$
\QM~\cite{Lopez2026}. The violations do not sit in a niche but in a broad connected region of
moderate memory and moderate wave amplitude; they are strongest in the periodic regime.

\subsection{The dynamic test}

The new step is a test in three stages \QM~\cite{Lopez2026}. First the droplets run with the coupling
switched on and the settings $a$, $a^*$ until they are synchronized. Then the coupling is switched
off ($\zeta=0$), and both subsystems run on independently for a time $T_2$. Finally the settings are
changed to $b$, $b^*$ by raising the barriers; this quickly drives each droplet into one of the two
wells, and that well is the outcome. The mean $\langle|S|\rangle$ oscillates with $T_2$ and exceeds
the bound 2 periodically; the maximum lies at about $2.85$ \QM. The authors summarize the result by
saying that \enquote{nonlinear deterministic dynamics can produce Bell violations through
wave-mediated synchronization mechanisms} \QM~\cite{Lopez2026}, and elsewhere more briefly:
\enquote{memory (specifically, time-delay coupling) can account for entanglement}
\QM~\cite{Lopez2026}.

What matters is how the authors themselves explain the violation. The hidden variables are not only
the position and velocity of the droplets but also the extended wave field, and its form depends on
the geometry of the wells. The violations are therefore to be understood in the sense that
\enquote{the wave form is influenced by the measurement settings} \QM~\cite{Lopez2026}. The
measurement-independence premise, $p(\lambda\,|\,a,b)=p(\lambda)$, is thus not satisfied in the
droplet system, and the authors say so explicitly. Vervoort had asked earlier whether this premise
can hold at all in systems with a background field \QM~\cite{Vervoort2018}.

% ============================================================
\section{Recomputing the bounds}
% ============================================================

For the classification three bounds on the CHSH value $S=E(a,b)+E(a',b)+E(a,b')-E(a',b')$ suffice:
the Bell bound $|S|\le2$ for local models with measurement independence \EM~\cite{Bell1964,CHSH1969},
the Tsirelson bound $|S|\le2\sqrt2=2.828$ for quantum correlations \EM~\cite{Tsirelson1980}, and
the algebraic bound $|S|\le4$, which follows from $|E|\le1$.

\begin{center}
\small
\begin{tabular}{@{}p{6.2cm}p{1.4cm}p{6.6cm}@{}}
\toprule
\textbf{Value} & \textbf{$S$} & \textbf{Position} \\
\midrule
Papatryfonos et al., static & $2.49$ & above 2, below $2\sqrt2$ \KM \\
López et al., dynamic (maximum) & $2.85$ & $0.022$ above $2\sqrt2$, i.e.\ $0.76\,\%$ \KM \\
López et al., static (mean) & $3.15$ & $0.32$ above $2\sqrt2$, about 214 standard errors \KM \\
\bottomrule
\end{tabular}
\end{center}

The values of the reduced model lie above the Tsirelson bound. The authors note this themselves for
the static value. A statement follows that matters for every deterministic reading: the droplet
correlations are not quantum correlations but stronger ones. This is not a flaw of the paper; it lies
in the mechanism. If the distribution of the hidden variables depends on the settings, any
correlation table can be produced. The check script shows this with a deterministic model that gives
unequal outcomes for the pair $(a',b')$ and equal ones otherwise; it reaches the algebraic bound $S=4$
\KM. Measurement dependence thus explains why the Bell bound 2 can fall. It does not explain why
nature stops at $2\sqrt2$.

A remark on the reach of the dynamic test. During the coupling phase the settings $a$, $a^*$ are
already fixed, and they shape the wave and the synchronized state that continues after the
switch-off. The later settings $b$, $b^*$ are rigidly paired with $a$, $a^*$. The test therefore
excludes signalling after the switch-off, but not the dependence of the hidden variables on the
settings; this agrees with the authors' reading. It is therefore not a counterpart of the
loophole-free Bell tests with randomly and late chosen settings \EM~\cite{Hensen2015}, but an example
of what measurement dependence concretely looks like in a classical field \SM.

% ============================================================
\section{What the papers show}
% ============================================================

First: a fully deterministic classical system with a real field can violate the Bell inequality,
without action at a distance and without randomness in the single run. The randomness of the
measurement series comes from the chaos of the trajectory and from the initial conditions, as with a
die; the authors draw this comparison themselves \QM~\cite{Lopez2026}. This is an existence statement,
and one model suffices for it.

Second: the price is measurement dependence. Without it the Bell bound holds for every local model,
and no mechanism can exceed it \EM. With it the bound is open up to 4. This way out is called
superdeterminism \QM~\cite{HossenfelderPalmer2020}. Palmer, who has advocated superdeterminism for
years, writes soberly: \enquote{No experiment to date has closed the superdeterministic
\enquote{loophole}} \QM~\cite{Palmer2023}. The droplets do not close this loophole; they give it a
tangible mechanics.

Third: the strength of the correlation is, in the droplet system, a matter of parameters, not of a
principle. The reduced model exceeds Tsirelson; the full model of Papatryfonos et al.\ stays below
it. A bound at $2\sqrt2$ arises from the structure of neither model.

% ============================================================
\section{Placing the result in the FFGFT}
% ============================================================

\subsection{Common ground}

Measurement in the FFGFT is a deterministic transition to a pole (Doc.~230): the outcome is
$A(z,\lambda)=\operatorname{sgn}(z-\lambda)$, fixed by the position $z$ of the system and the
quantity $\lambda$ of the measuring arrangement, and this fixing happens in the interaction between
system and measuring arrangement. This is the same logical point at which the droplets leave
measurement independence: the setting shapes the quantity that determines the outcome. Doc.~230
therefore calls the model contextual, and the Kochen--Specker theorem does not apply. The droplets
thus provide a concrete classical example of the kind of contextuality the FFGFT assumes \SM.

The reading of randomness agrees as well. In both cases the single run is fixed, and the frequencies
of the measurement series arise because the determining quantity is not controlled. In Doc.~230 the
ignorance concerns only $\lambda$; with the droplets it concerns the initial conditions and the wave
field. And finally, measurement is coupling in both cases: Doc.~328 describes every measurement as
the coupling of a readout system with a degree of coupling $\kappa$, and the droplets are measured by
raising a barrier until the system falls into a well.

\subsection{Differences}

\textbf{Coupling strength.} In the droplet model the dimensionless wave amplitude is of order 5 at a
damping of 1; the coupling is strong, and synchronization is the result. In the FFGFT the effective
coupling in the circle map is $K_{\mathrm{eff}}=2\pi\xi\approx8.4\times10^{-4}$, and the Arnold tongue
of order $1/2$ is $5.6\times10^{-8}$ wide (Doc.~328) \KM. The ratio is about 6000 \KM. Dynamic
synchronization of the kind the droplets carry is excluded as the mechanism of the correlation in
the FFGFT at this coupling \XM. This agrees with Docs.~328 and 343: the discreteness of the FFGFT is
topological, not dynamic. The Bell connection in the FFGFT is a topologically connected cylinder
configuration on $T^4$ (Doc.~230), not a locked state of two oscillators.

\textbf{The Tsirelson bound.} Here lies the most important difference. For the singlet, the quotients
from Doc.~230 give the fraction of equal outcomes $\sin^2\!\big((a-b)/2\big)$ and from it
$E(a,b)=-\cos(a-b)$ \KM. This correlation reaches exactly $2\sqrt2$ at the angles $0$, $90^\circ$,
$45^\circ$, $-45^\circ$ and exceeds it for no choice of angles; a grid search in the check script
confirms this \KM. The bound is here not an additional assumption but the maximum of a combination
of cosines. Where the droplets go beyond $2\sqrt2$ because measurement dependence alone limits
nothing, the FFGFT stays at the bound because its weights are geometric quotients (Docs.~230, 367,
370).

\textbf{What remains open.} The local model $A=\operatorname{sgn}(z-\lambda)$ with a common $\lambda$
independent of the settings gives for two subsystems only the sawtooth correlation
$E=-1+2|a-b|/\pi$ and thus at most $S=2$ \KM. The cosine form requires more: the topological
connection on $T^4$ that does not let the two cylinder configurations be independent. In the corpus
the agreement with $-\cos(a-b)$ is carried by the complete translatability between Hilbert space and
cylinder representation (Doc.~230); a derivation of the cosine form from the topological connection
itself, without recourse to Hilbert space, is not available \SM. This is the point at which the
droplets set the FFGFT a precise task: to show that the connection on $T^4$ produces exactly the
cosine form and nothing more.

\textbf{Memory and time.} The droplets need a memory of finite duration, and their correlation
oscillates with the waiting time $T_2$. In the FFGFT each mode carries its own time scale
$\tilde T=1/m$ (Doc.~334); the correlation is not a stored history but a property of the
configuration. The FFGFT therefore does not expect the Bell correlation to oscillate with the waiting
time between creation and measurement \SM.

\subsection{What the FFGFT takes from this}

The droplets sharpen the question by which a deterministic reading must be measured. That Bell
violations are possible without action at a distance and without randomness is shown vividly by
them. The real question is why nature stops at exactly $2\sqrt2$. A reading that offers only
measurement dependence does not answer it. The FFGFT has an answer in outline: the Born weights as
geometric quotients of a closed geometry. The answer is worked out for the single case (Docs.~230,
367); for the entangled pair it is still outstanding \SM.

% ============================================================
\section{What does not match}
% ============================================================

\begin{itemize}[leftmargin=2em,itemsep=3pt,topsep=3pt]
  \item The droplets live in one dimension and in a model with an idealized wave kernel; the authors
        themselves name as open whether the violations persist in a fully dynamic, local wave model
        with finite propagation time \QM~\cite{Lopez2026}.
  \item The results are simulations; an experiment with real droplets does not yet exist.
  \item The droplets' measurement setting is a geometry that shapes wave and trajectory; in the FFGFT
        $\lambda$ is a quantity of the measuring arrangement on the same axis as $z$, without an
        extended field in position space.
  \item The droplets exceed Tsirelson, the FFGFT does not.
  \item The droplets synchronize dynamically at strong coupling; the FFGFT is deeply subcritical at
        $K_{\mathrm{eff}}=2\pi\xi$.
\end{itemize}

% ============================================================
\section{Overview}
% ============================================================
\begin{center}
\small
\begin{tabular}{@{}p{8.2cm}p{1.2cm}p{4.4cm}@{}}
\toprule
\textbf{Statement} & \textbf{Status} & \textbf{Basis} \\
\midrule
Droplets violate CHSH: $2.49$; $3.15$; dynamic up to $2.85$ & \QM & \cite{Papatryfonos2023,Lopez2026} \\
Explanation via measurement dependence of the wave & \QM & \cite{Lopez2026} \\
$3.15$ and $2.85$ lie above $2\sqrt2$ & \KM & check script \\
Measurement dependence allows $S$ up to 4 & \KM & check script \\
Born quotients from Doc.~230 give $-\cos(a-b)$, maximum exactly $2\sqrt2$ & \KM & check script \\
$\operatorname{sgn}(z-\lambda)$ with common $\lambda$: sawtooth, $S\le2$ & \KM & check script \\
Contextual fixing in Doc.~230 corresponds to the droplets' measurement dependence & \SM & classification \\
Synchronization as FFGFT mechanism at $K_{\mathrm{eff}}=2\pi\xi$ & \XM & Doc.~328, check script \\
Cosine form derived from the topological connection on $T^4$ & \SM & open \\
No oscillation of the correlation with the waiting time & \SM & Doc.~334 \\
\bottomrule
\end{tabular}
\end{center}

\begingroup\renewcommand{\chapter}[2]{}\small
\begin{thebibliography}{99}
\bibitem{Lopez2026} Á.~G.~López, R.~N.~Valani, Y.~Li, J.~W.~M.~Bush, \textit{Synchronization induces Bell violations in a model of walking droplets}, arXiv:2608.21915 (2026).
\bibitem{Papatryfonos2023} K.~Papatryfonos, L.~Vervoort, A.~Nachbin, M.~Labousse, J.~W.~M.~Bush, \textit{A platform for investigating Bell correlations in pilot-wave hydrodynamics}, arXiv:2208.08940v2 (2023).
\bibitem{Couder2005} Y.~Couder, S.~Protière, E.~Fort, A.~Boudaoud, \textit{Walking and orbiting droplets}, Nature \textbf{437} (2005) 208.
\bibitem{BushOza2021} J.~W.~M.~Bush, A.~U.~Oza, \textit{Hydrodynamic quantum analogs}, Reports on Progress in Physics \textbf{84} (2021) 017001.
\bibitem{Vervoort2018} L.~Vervoort, \textit{Are hidden-variable theories for pilot-wave systems possible?}, Foundations of Physics \textbf{48} (2018) 803.
\bibitem{Bell1964} J.~S.~Bell, \textit{On the Einstein Podolsky Rosen paradox}, Physics \textbf{1} (1964) 195.
\bibitem{CHSH1969} J.~F.~Clauser, M.~A.~Horne, A.~Shimony, R.~A.~Holt, \textit{Proposed experiment to test local hidden-variable theories}, Physical Review Letters \textbf{23} (1969) 880.
\bibitem{Tsirelson1980} B.~S.~Cirel'son, \textit{Quantum generalizations of Bell's inequality}, Letters in Mathematical Physics \textbf{4} (1980) 93.
\bibitem{Hensen2015} B.~Hensen et al., \textit{Loophole-free Bell inequality violation using electron spins separated by 1.3 kilometres}, Nature \textbf{526} (2015) 682.
\bibitem{HossenfelderPalmer2020} S.~Hossenfelder, T.~N.~Palmer, \textit{Rethinking superdeterminism}, Frontiers in Physics \textbf{8} (2020) 139.
\bibitem{Palmer2023} T.~N.~Palmer, \textit{Superdeterministic loophole}, letter, Physics Today, June 2023.
\end{thebibliography}
\endgroup
